% Style: ../template/dcptalk.sty (found through TEXINPUTS, see Makefile)
\usepackage[fonts = paper, codedir = gen]{dcptalk}
\input{gen/pygments}
\graphicspath{{images/}{../template/images/}}

% Colour names used by the diagrams
\colorlet{accent}{blue!60!black}
\colorlet{accent2}{orange!80!black}
\colorlet{bad}{dcpbad}
\colorlet{good}{dcpgood}

% Links (\url, \href) in the accent colour, so they stand out
\hypersetup{colorlinks = true, urlcolor = accent, linkcolor = black}

% Recursion schemes
% Latin Modern Math has no banana brackets: built as in the paper's lensbans
\newcommand\cata[1]{\mathopen{(\mkern-3.5mu|}#1\mathclose{|\mkern-3.5mu)}}
\newcommand\ana[1]{\mathopen{[\mkern-2mu(}#1\mathclose{)\mkern-2mu]}}
\newcommand\hylo[2]{\lBrack #1, #2 \rBrack}
\newcommand\ext[1]{\lBrack #1 \rBrack}
\newcommand\op[1]{\mathsf{#1}}
\newcommand\List{\op{List}}
\newcommand\inn{\op{in}}
\newcommand\out{\op{out}}
\usetikzlibrary{backgrounds}

% Under a code block that is not verbatim tool output or library source
\newcommand\simplified{\par{\hfill\footnotesize\textcolor{gray}{(simplified)}\par}}

% Code in mainstream languages (Python, SQL, Java, C) is orange, Haskell
% violet; dcptalk has Rocq blue, OCaml green, errors red, output grey
\newtcolorbox{mainbox}[1][]{enhanced, drop fuzzy shadow southeast,
  boxrule = 0pt, colframe = orange!12, colback = orange!12,
  halign = flush left, #1}
\newtcolorbox{hsbox}[1][]{enhanced, drop fuzzy shadow southeast,
  boxrule = 0pt, colframe = violet!8, colback = violet!8,
  halign = flush left, #1}
\makeatletter
\NewDocumentCommand\mainblock{O{} m}{\dcp@code{mainbox}{#1}{#2}}
\NewDocumentCommand\hsblock{O{} m}{\dcp@code{hsbox}{#1}{#2}}
\makeatother

% A tag in the top-right corner: which work a frame belongs to. At the end
% of the frame, like \slideref.
\NewDocumentCommand\worktag{O{accent} m}{%
  \begin{tikzpicture}[overlay, remember picture, alt = {Tag: #2}]
    \node[anchor = north east, draw = #1, fill = #1!12, text = #1,
      rounded corners = 2pt, font = \footnotesize\bfseries, inner sep = 3pt]
      at ($(current page.north east) + (-4mm, -4mm)$) {#2};
  \end{tikzpicture}\par}
\newcommand\rocqtag{\worktag[accent]{Rocq Library}}
\newcommand\icfptag{\worktag[accent]{ICFP 2016}}
\newcommand\ongoingtag{\worktag[accent2]{ongoing}}

% Slides that the live demo covers (skip them after a good demo): a DEMO
% tag in the top-right corner, in place of the work tag. At the end of the
% frame.
\newcommand\demotag{\worktag[bad]{DEMO}}

% \obstacle{text}: a small note under the work tag, naming which of the
% "Three Obstacles" (termination, no extra axioms, clean code) a frame is
% about. At the end of the frame; a zero-height box, as for \callout, so the
% frame's content does not move.
\makeatletter
\NewDocumentCommand\obstacle{m}{%
  \par\edef\dcp@prevdepth{\the\prevdepth}\nointerlineskip\vbox to 0pt{%
    \begin{tikzpicture}[overlay, remember picture, alt = {Obstacle: #1}]
      \node[anchor = north east, draw = golden!80!black, fill = yellow!20,
        rounded corners = 2pt, font = \scriptsize, inner sep = 2pt]
        at ($(current page.north east) + (-4mm, -10.5mm)$) {#1};
    \end{tikzpicture}\vss}%
  \prevdepth = \dcp@prevdepth\relax}
\makeatother
% \openquestion{text}: the same chip, naming which open question (layout,
% scheduling, search) a low-level frame is about. At the end of the frame.
\NewDocumentCommand\openquestion{m}{\obstacle{open question: #1}}

\usetikzlibrary{shadows}

% Inside \slideref: the next reference starts on a new line (no \par: the
% text also goes into the picture's alt text)
\newcommand\rfsep{\\}
% This deck's references sit a little lower than dcptalk's (4mm from the
% bottom, not 7mm), to make room for one reference per line, under a short
% rule that separates them from the slide's content (as for footnotes)
\RenewDocumentCommand\slideref{+m}{%
  \begin{tikzpicture}[overlay, remember picture, alt = {Reference: #1}]
    \node[anchor = south west, font = \tiny, text width = 150mm,
      align = left] at ($(current page.south west) + (10mm, 4mm)$)
      {\textcolor{black!40}{\rule{25mm}{0.4pt}}\\[0.3ex] #1};
  \end{tikzpicture}\par}

% A short citation on its own line, in small grey type (timeline boxes)
\newcommand\tlcite[1]{\\{\scriptsize\textcolor{gray}{#1}}}

% \excerpt<overlay>{quotation}{source}: a quotation from a paper in a
% sticky note over the slide, revealed on the given overlay. At the end of
% the frame.
\NewDocumentCommand\excerpt{D<>{2-} +m +m}{%
  \only<#1>{%
    \begin{tikzpicture}[overlay, remember picture, alt = {Quotation: #2 (#3)}]
      \node[anchor = north, fill = yellow!20, draw = golden!80!black,
        rounded corners = 3pt, drop shadow, text width = 140mm,
        align = justify, font = \small, inner sep = 4mm]
        at ($(current page.north) + (0, -18mm)$)
        {#2\par\smallskip{\raggedleft\footnotesize --- #3\par}};
    \end{tikzpicture}}\par}

% Footer: "frame/total", bottom right. ltx-talk counts frames, not overlays.
% (\hfill in font pushes it right: the footer has no alignment key.)
\EditInstance{footer}{std}{
  color         = black!50,
  font          = \hfill\scriptsize,
  element-order = {framenumber, totalframes},
  separator     = /,
  right-hspace  = 4mm,
}

% The title and Part frames have no footer, and are not counted. (A counter
% of our own: ltx-talk's frame counter also sets the PDF page labels, which
% pdfpc uses to group overlays.)
\makeatletter
\newcounter{dcpskipped}
\def\ps@talknofooter{\ps@talk\def\@oddfoot{}\let\@evenfoot\@oddfoot}
\NewCommandCopy\dcp@titleframe\titleframe
\RenewDocumentCommand\titleframe{+m +m +m}{\stepcounter{dcpskipped}%
  \dcp@titleframe{#1}{#2}{#3\thispagestyle{talknofooter}}}
\NewCommandCopy\dcp@partframe\partframe
\RenewDocumentCommand\partframe{m m}{\stepcounter{dcpskipped}%
  \dcp@partframe{#1}{#2\thispagestyle{talknofooter}}}
\ExplSyntaxOn
\property_new:nnnn { dcpskipped } { shipout } { 0 } { \int_use:N \c@dcpskipped }
\AddToHook { enddocument / afterlastpage }
  { \property_record:nn { dcp-lastpage } { dcpskipped } }
% Backup frames (after \appendix): B1, B2, ... (the first appendix frame is
% the Backup divider); the total counts the talk's frames only
\cs_set:Npn \@framenumber
  {
    \bool_if:NTF \l__talk_appendix_bool
      { B \int_eval:n { \c@appendixframe - 1 } }
      { \int_eval:n { \c@frame - \c@dcpskipped } }
  }
\cs_set:Npn \@totalframes
  {
    \int_eval:n
      {
          \property_ref:nn { lastpage } { frames }
        - \property_ref:nn { dcp-lastpage } { dcpskipped }
      }
  }
\ExplSyntaxOff
% \backuppart, after "Thank you!": an uncounted divider, then the backup
% frames, numbered without a total
\NewDocumentCommand\backuppart{}{%
  \appendix
  \EditInstance{footer}{std}{element-order = {framenumber}}%
  \dcp@partframe{Backup}{}}
\makeatother
